\section{A fixed surface bundle with arbitrarily small pointwise Einstein defect}
\label{sec:fixed-bundle-geometry}

The distinction between integral separation and unrestricted Einstein
existence can be tested on a fixed smooth manifold.  The examples in this
section have actual monopole classes in both orientations and admit
pointwise Einstein approximation with bounded curvature derivatives.
Their topology excludes complex and nonpositively curved metrics, but
does not exclude all negative Einstein metrics.

\subsection{The smooth object and its topological invariants}

Fix integers $p,q\ge2$, oriented surfaces $F=\Sigma_p$ and $B=\Sigma_q$,
a nonseparating simple closed curve $\alpha\subset F$, and a positive
Dehn twist $D$ supported in an annular neighborhood of $\alpha$.
Choose a base point outside this annulus and a primitive epimorphism
$u:\pi_1(B)\twoheadrightarrow\mathbb Z$.  Define
\begin{equation}\label{eq:fixed-bundle-object}
 E_D=(\widetilde B\times F)/\pi_1(B),\qquad
 \gamma\cdot(x,y)=(\gamma x,D^{u(\gamma)}y).
\end{equation}
This is an oriented smooth surface bundle over $B$.  The twist, the
epimorphism, the genera, and the smooth total space will remain fixed
throughout the metric construction.  Put
\begin{equation}\label{eq:fixed-bundle-volume}
 c=\chi(E_D)=4(p-1)(q-1),\qquad
 V_0=4\pi^2c=16\pi^2(p-1)(q-1).
\end{equation}

\begin{proposition}\label{prop:fixed-bundle-topology}
The manifold $E_D$ is aspherical, and
\begin{align}
 \pi_1(E_D)&=\Pi_p\rtimes_{D_*^{u(\cdot)}}\Pi_q,
       &\Pi_r&=\pi_1(\Sigma_r),\label{eq:fixed-bundle-group}\\
 \tau(E_D)&=0,
       &H_1(E_D;\mathbb Z)&=\mathbb Z^{2p+2q-1}.
       \label{eq:fixed-bundle-invariants}
\end{align}
In particular $b_2^+(E_D)=b_2^-(E_D)=2pq$.  It has symplectic
structures in both orientations, with canonical classes of square $2c$
in their respective orientations.
\end{proposition}

\begin{proof}
The chosen fixed point of $D$ gives a section of the bundle, and hence
the splitting in \eqref{eq:fixed-bundle-group}.  The homotopy sequence
of the fibration, together with asphericity of the two surfaces, shows
that all higher homotopy groups of $E_D$ vanish.  Euler characteristic
is multiplicative, giving \eqref{eq:fixed-bundle-volume}.

The bundle is also the pullback of the mapping torus of $D$ along a
smooth map $B\to S^1$ representing $u$.  If $V$ denotes its vertical
tangent bundle, $e(V)^2$ is therefore pulled back from a three-manifold
and vanishes.  With $e_B=\pi^*e(TB)$ we have $e_B^2=0$ and
\[
 p_1(TE_D)=e(V)^2+e_B^2=0.
\]
The signature theorem proves $\tau(E_D)=0$.

Abelianizing the split extension gives
\[
 H_1(E_D;\mathbb Z)
 =H_1(B;\mathbb Z)\oplus
   H_1(F;\mathbb Z)/(D_*-I)H_1(F;\mathbb Z).
\]
The homological action of a nonseparating twist is the transvection
$x\mapsto x\pm(x\cdot[\alpha])[\alpha]$.  Since $[\alpha]$ is
primitive, the image of $D_*-I$ is the primitive rank-one subgroup
generated by $[\alpha]$.  This proves the asserted integral homology.
The Betti-number formula now follows from $\chi=2-2b_1+b_2$ and
$\tau=0$.

Choose an area form $\mu_F$ preserved by the annular twist, and a
positive area form $\mu_B$ on the base.  On the flat-bundle presentation
\eqref{eq:fixed-bundle-object}, the forms
$\Omega_+=\mu_F+\pi^*\mu_B$ and
$\Omega_-=\mu_F-\pi^*\mu_B$ descend to closed nondegenerate forms.
Their squares have opposite signs.  The canonical classes can be taken
as $-e(V)-e_B$ and $-e(V)+e_B$, respectively.  Because
$\langle e(V)e_B,[E_D]\rangle=c$, their squares are $2c$ in the
corresponding symplectic orientations.
\end{proof}

The same construction has a useful description in the notation of
Sections~\ref{sec:closed-shear} and~\ref{sec:shear-topology}.
Cutting $F$ along $\alpha$ gives $F_{p-1,2}$.  Then $E_D$ is the
self-gluing of $F_{p-1,2}\times\Sigma_q$ with identity mapping class
on $\Sigma_q$ and primitive circle-valued shear representing $u$.
Indeed, on a simply connected base chart a real lift of this shear
can be absorbed by rotation in the cut annulus.  Around a base loop
the lift changes by $u(\gamma)$ full turns, whose return map is
$D^{u(\gamma)}$.  Thus this self-gluing gives precisely
\eqref{eq:fixed-bundle-object}, not a different smooth candidate.
In particular, the virtual K\"ahler obstruction of
Section~\ref{sec:virtual-nonkahler} applies for every $p,q\ge2$.

\subsection{A closed slow-interpolation metric}

For a metric $G$ on a four-manifold write
\[
 \Phi(G)=\Ric_G+G,\qquad
 \delta_-(G)=\sup\bigl|G^{-1}\Phi(G)\bigr|_{\mathrm{op}}.
\]
The endomorphism norm in this definition is distinct from the tensor
norm used in integrated curvature energies.

\begin{theorem}\label{thm:fixed-bundle-approximation}
For every choice of the fixed data in \eqref{eq:fixed-bundle-object},
there are smooth metrics $G_\ell$, $\ell\downarrow0$, on that same
$E_D$ such that
\begin{gather}
 \Vol(E_D,G_\ell)=V_0,\qquad \delta_-(G_\ell)\longrightarrow0,
 \label{eq:fixed-bundle-approximation}\\
 \|\nabla^m\Rm_{G_\ell}\|_\infty\le C_m,\qquad
 \|\nabla^m\Phi(G_\ell)\|_\infty\le C_m/L_\ell
       \quad(m\ge0),\label{eq:fixed-bundle-derivatives}\\
 L_\ell=\frac12\operatorname{arsinh}\frac1{\sinh(\ell/2)}
       \longrightarrow\infty.\nonumber
\end{gather}
The constants may depend on the fixed surfaces, twist profile, and
derivative order.  The injectivity radius tends to zero and the
diameter tends to infinity.  Every connected $d$-sheeted finite cover
has the pulled-back family with the same pointwise estimates and
volume $dV_0$.
\end{theorem}

\begin{proof}
An isotopy of the annular representative of $D$ induces an isomorphism
of its mapping torus and hence of the pullback bundle.  We may therefore
use the following representative, and pull the resulting metrics back
to the original fixed smooth bundle.
Choose a hyperbolic metric $h_0$ on $F$ for which $\alpha$ is geodesic.
On a fixed subcollar, with circle coordinate $s\in\mathbb R/\mathbb Z$,
write
\[
 h_0=dr^2+a(r)^2ds^2,\qquad a(r)=a_0\cosh r.
\]
Choose a smooth $\zeta$ equal to $0$ and $1$ near its two boundary
components, so $D(r,s)=(r,s+\zeta(r))$.  The integer displacement
at the second boundary makes $D$ the identity there.  Define
\begin{equation}\label{eq:fixed-bundle-fiber-path}
 h_z=dr^2+a(r)^2\bigl(ds+z\zeta'(r)\,dr\bigr)^2,
 \qquad 0\le z\le1,
\end{equation}
and take $h_z=h_0$ off this collar.  The local substitution
$s'=s+z\zeta(r)$ proves
\[
 K_{h_z}=-1,\qquad d\mu_{h_z}=d\mu_{h_0},
 \qquad h_1=D^*h_0.
\]
For nonintegral $z$ this local substitution need not extend to a
diffeomorphism of the closed fiber.  Only the metrics are patched
across the collar boundaries, where $\zeta'=0$.

Represent the Poincar\'e dual of the primitive class $u$ by a
nonseparating base curve.  Choose hyperbolic metrics on $B$ for which
its length is $\ell$ and the other Fenchel--Nielsen data are fixed.
The collar theorem supplies
\begin{equation}\label{eq:fixed-bundle-base-collar}
 dt^2+b(t)^2d\theta^2,\qquad b(t)=\ell\cosh t,\qquad
 |t|\le L_\ell,\quad\theta\in\mathbb R/\mathbb Z;
\end{equation}
we use half the available collar width
\cite[Appendix A]{RupflinToppingCollars}.  Fix a smooth function
$\eta:[-1,1]\to[0,1]$, constant $0$ near $-1$ and $1$ near $1$,
and set $z(t)=\eta(t/L_\ell)$.  On this base collar put
\begin{equation}\label{eq:fixed-bundle-metric}
 G_\ell=dt^2+b(t)^2d\theta^2+h_{z(t)}.
\end{equation}
On its complement use the product of the base metric with $h_0$.
The two ends are glued by the identity and by $D$, respectively.
At the second end $h_1=D^*h_0$, and the profile is constant near
each end.  Consequently every transverse derivative matches.  This
is a positive smooth metric on the fixed bundle, not just a metric
on its cut-open part.

We compute all Ricci components.  Put
\[
 j=\tanh r,\qquad k=\tanh t,\qquad
 Q=a\zeta',\qquad v=Qz'.
\]
The orthonormal coframe
\[
 (e^0,e^1,e^2,e^3)
   =\bigl(dt,b\,d\theta,dr,a(ds+z\zeta'\,dr)\bigr)
\]
satisfies
\begin{equation}\label{eq:fixed-bundle-structure}
 de^1=k\,e^0\wedge e^1,\qquad
 de^3=j\,e^2\wedge e^3+v\,e^0\wedge e^2,\qquad
 de^0=de^2=0.
\end{equation}
On the fiber, the shape endomorphism and its covariant derivative
along $t$ are
\begin{equation}\label{eq:fixed-bundle-shape}
 S=\frac12\begin{pmatrix}0&v\\v&0\end{pmatrix},
 \quad D_tS=\begin{pmatrix}v^2/2&v_t/2\\v_t/2&-v^2/2\end{pmatrix},
 \quad \operatorname{div}_{h_z}S=\frac{v_r+2jv}{2}e^3.
\end{equation}
For example, the time connection in this frame is
$\nabla_{e_0}e_2=-(v/2)e_3$ and
$\nabla_{e_0}e_3=(v/2)e_2$.  Its commutator with the matrix of
$S$ gives the quadratic diagonal terms in $D_tS$; these terms
would be lost by differentiating the matrix entries alone.
Since $\operatorname{tr}S=0$, the slice Gauss--Codazzi equations read
\[
 \Phi_{tt}=-\operatorname{tr}S^2,\quad
 \Phi_{11}=0,\quad
 \Phi_{FF}^{\sharp}=-D_tS-kS,\quad
 \Phi_{tF}=\operatorname{div}_{h_z}S.
\]
They give the complete defect matrix
\begin{equation}\label{eq:fixed-bundle-ricci-matrix}
 \Phi(G_\ell)_{\mathrm{ON}}=
 \begin{pmatrix}
 -v^2/2&0&0&(v_r+2jv)/2\\
 0&0&0&0\\
 0&0&-v^2/2&-(v_t+kv)/2\\
 (v_r+2jv)/2&0&-(v_t+kv)/2&v^2/2
 \end{pmatrix}.
\end{equation}
Thus, with
$A_0=\|Q\|_\infty$, $A_1=\|Q'+2jQ\|_\infty$, and
$M_i=\|\eta^{(i)}\|_\infty$, the diagonal and off-diagonal parts give
\begin{equation}\label{eq:fixed-bundle-explicit-bound}
 \delta_-(G_\ell)
 \le\frac{A_0^2M_1^2}{2L_\ell^2}
 +\frac1{2L_\ell}
   \sqrt{A_1^2M_1^2+A_0^2(M_1+M_2/L_\ell)^2}.
\end{equation}
The off-diagonal part is a single row and column with two nonzero
entries, whose operator norm is their Euclidean length.  This
explains the square root without discarding either mixed Ricci block.
The trace and volume form yield the exact identities
\begin{equation}\label{eq:fixed-bundle-scalar-volume}
 s_{G_\ell}=-4-\frac12v^2,\qquad
 d\mu_{G_\ell}=d\mu_{B,\ell}\,d\mu_{h_0},\qquad
 \Vol(G_\ell)=V_0.
\end{equation}

There is no inverse-$\ell$ loss in differentiating the curvature.
All coefficients are independent of $\theta$, and the circle
connection uses $b'/b=\tanh t$, whose derivatives are bounded.
The fiber geometry and path are fixed, and each derivative of $z$
contributes its corresponding inverse power of $L_\ell$.
The structure equations \eqref{eq:fixed-bundle-structure}, or a fixed
finite atlas outside the fiber collar, therefore give
\eqref{eq:fixed-bundle-derivatives} at every fixed order.  Outside the
active region the metric is an exact product.

The base waist is noncontractible; the section through the fixed
fiber point lifts it to a loop of length $\ell$.  Hence
$\operatorname{inj}_{\min}(G_\ell)\le\ell/2$.  Projection to the base
does not increase lengths, and the base collar has diverging width,
so the diameter diverges.  The region over $|t|\le L_\ell$ has volume
\[
 8\pi(p-1)\ell\sinh L_\ell=O(\sqrt\ell).
\]
The thick-part limit is
$\Sigma_p\times\Sigma_{q-1,2}$, with the complete finite-area
hyperbolic metric on the second factor and total volume $V_0$.
Finally, a finite covering is a local isometry for each pulled-back
metric; it preserves every pointwise estimate and multiplies volume
by its degree.
\end{proof}

\begin{remark}\label{rem:fixed-bundle-pointwise}
The small defect in \eqref{eq:fixed-bundle-explicit-bound} is a global
maximum estimate; the shrinking volume of the active region is not
its justification.  The same computation gives
$K(e_0,e_3)=v^2/4>0$ wherever $v\ne0$.  Moreover no $G_\ell$ is
Einstein: its $11$ Ricci entry forces the Einstein constant to be
$-1$, after which its $00$ entry forces $v=0$, contrary to the
nontrivial twist.  This last argument concerns the displayed metrics,
not arbitrary metrics on $E_D$.
\end{remark}

\subsection{Finite covers and the topological restrictions}

\begin{proposition}\label{prop:fixed-bundle-cover-obstructions}
Every connected finite cover of $E_D$ is a surface bundle over a
surface with both genera at least two and with infinite monodromy.
No such cover admits an integrable complex structure, in either
orientation, or a metric of nonpositive sectional curvature.
\end{proposition}

\begin{proof}
Let $G=\pi_1(E_D)$, $N=\Pi_p$, and let $H<G$ have finite index.
Set $K=H\cap N$ and let $Q'<\Pi_q$ be the image of $H$.  If
$a=[N:K]$ and $b=[\Pi_q:Q']$, the degree is $d=ab$.  The bundle
projection lifts to the base cover corresponding to $Q'$.  Its
proper submersion has connected fibers: the component cover of a
proper submersion is determined by the image on fundamental groups,
and that image is all of $Q'$.  Equivalently, this follows by applying
the connected-fiber factorization to the lifted projection.  The
fiber and base genera are consequently
\begin{equation}\label{eq:fixed-bundle-cover-genera}
 p'=1+a(p-1),\qquad q'=1+b(q-1).
\end{equation}
The signature of the total space is still zero.

Choose $h\in H$ whose base image has $u$-value $n\ne0$; this is
possible because a nonzero homomorphism to $\mathbb Z$ remains
nonzero on every finite-index subgroup.  Conjugation by $h$ defines
an automorphism $\alpha_h$ of $N$ preserving $K$, with outer class
$[D]^n$.  If the covering monodromy were finite, there would be
$m\ge1$ and $k_0\in K$ such that
$\alpha_h^m|_K=\operatorname{Inn}(k_0)|_K$.
Let $\beta=\operatorname{Inn}(k_0^{-1})\alpha_h^m$.
For every $x\in N$ some positive power $x^r$ lies in $K$, whence
$\beta(x)^r=x^r$.  Surface groups have unique roots, so $\beta(x)=x$.
For completeness, unique roots follow from their torsion-free cyclic
centralizers: two elements with the same nonidentity power lie in
the same cyclic centralizer and have the same exponent.  It follows
that $\alpha_h^m$ is inner on all of $N$, contradicting the nontrivial
transvection $D_*^{nm}$ on $H_1(F;\mathbb Z)$.
Kotschick's characterization of zero-signature surface bundles
\cite[Theorem 4]{KotschickFibered} now excludes a complex structure.
This theorem concerns any complex structure on the total space; it
does not assume that the displayed projection is holomorphic.

Here is the nonpositive-curvature obstruction in group-theoretic
form.  The primitive monodromy gives a subgroup
$\Pi_p\rtimes_{D_*}\mathbb Z<G$.  Choose surface generators so
that its relations include
\[
 zaz^{-1}=a,\quad zbz^{-1}=ba,\quad
 zc_iz^{-1}=c_i,\quad zd_iz^{-1}=d_i,\qquad
 [a,b]\prod_{i=2}^p[c_i,d_i]=1.
\]
Set $v=bab^{-1}$.  The elements $a,v,c_i,d_i$ commute with $z$.
Suppose this group acted properly by semisimple isometries on a
complete CAT$(0)$ space.  The infinite-order elements are axial,
and the splitting
$\operatorname{Min}(z)=Y\times\mathbb R$ gives a homomorphism
$\tau_z:C(z)\to\mathbb R$ measuring translation in the oriented
$z$-direction.  The surface relation gives
$\tau_z(a)=\tau_z(v)$.
On the other hand, $az=b^{-1}zb$ and $v^{-1}z=bzb^{-1}$.
The flat-torus theorem, applied separately to the commuting pairs
$(a,z)$ and $(v,z)$, therefore gives
\[
 \tau_z(a)=-\frac{\ell(a)^2}{2\ell(z)}<0,\qquad
 \tau_z(v)=\frac{\ell(v)^2}{2\ell(z)}>0.
\]
Here $\ell$ denotes translation length.  Both formulas use the same
translation homomorphism of $C(z)$, even though the two flats need
not coincide.  This is a contradiction.

Proper semisimple CAT$(0)$ actions restrict to subgroups, and their
existence is invariant under passage between finite-index groups
\cite[Theorem 2.3 and the proof of Theorem 4.5]{KapovichLeeb}.
Thus no finite-index subgroup of $G$ has such an action.  A closed
nonpositively curved metric would provide one by deck transformations
on its universal cover.  This excludes that metric on every cover.
\end{proof}

The stronger assertion that no finite-index subgroup of $G$ is a
K\"ahler group follows from the self-gluing identification and
Section~\ref{sec:virtual-nonkahler}.  It is not an inference from
the absence of complex structures on these particular four-manifolds.
Also, the finite-cover argument does not claim that every cover has
odd first Betti number.

\subsection{Actual monopoles and the allowed Einstein geometry}

Let $\widehat E\to E_D$ be a connected $d$-sheeted cover and write
$V_*=dV_0$ and $\widehat c=\chi(\widehat E)=dc$.
The pulled-back symplectic forms give canonical classes of square
$2\widehat c>0$ in each orientation.  The fibration
\eqref{eq:fixed-bundle-cover-genera} gives $b_1(\widehat E)\ge2q'$,
so $b_2^\pm=\widehat c/2-1+b_1>1$.
Taubes's theorem \cite{Taubes1994} therefore supplies actual
unperturbed monopoles for every metric, separately in the two
orientations.  Since their self-dual projection square is at least
$2\widehat c$, reducibility cannot occur in the limit from generic small
perturbations.

In particular, the scalar negative-part estimate gives
\begin{equation}\label{eq:fixed-bundle-scalar-budget}
 \int_{\widehat E}s_-^2\,d\mu\ge64\pi^2\widehat c=16V_*.
\end{equation}
Equality forces the nonzero monopole spinor to be parallel and hence
the metric to be K\"ahler \cite{LeBrun2001}.  Proposition
\ref{prop:fixed-bundle-cover-obstructions} excludes equality here.
The pulled-back family has $\E(G_\ell)\le4\delta_-(G_\ell)^2\to0$.
These facts coexist because the canonical square is $2\chi$:
the actual class has topological ratio $T/a^2=1$, not a critical
ratio required for a topology-driven positive separation.

\begin{proposition}\label{prop:fixed-bundle-einstein-window}
Any Einstein metric on $\widehat E$ has negative Einstein constant.
After normalization to $\Ric_g=-g$ it satisfies
\begin{equation}\label{eq:fixed-bundle-einstein-window}
 V_*<\Vol_g(\widehat E)<3V_*.
\end{equation}
It has a positive-sectional-curvature plane.  Each Weyl half has
three distinct eigenvalues on a common open dense subset of
$\widehat E$.
\end{proposition}

\begin{proof}
Nonnegative scalar curvature contradicts
\eqref{eq:fixed-bundle-scalar-budget}.  At $\Ric=-g$ the same
inequality gives $16\Vol(g)\ge16V_*$, strictly because equality
would be K\"ahler.  The Einstein Gauss--Bonnet identity, with
operator norms on the Weyl halves, gives
\begin{equation}\label{eq:fixed-bundle-einstein-cgb}
 2V_*=8\pi^2\widehat c
   =\int_{\widehat E}(|W^+|^2+|W^-|^2)\,d\mu
      +\frac23\Vol(g).
\end{equation}
It follows that $\Vol(g)\le3V_*$.  Equality would force $W=0$
and hence sectional curvature $-1/3$, contrary to Proposition
\ref{prop:fixed-bundle-cover-obstructions}.  The same proposition
requires some plane of positive sectional curvature.

Neither Weyl half is identically zero: since $\tau=0$, the
signature formula would force the other half to vanish as well.
If, say, $W^+$ had at most two eigenvalues everywhere,
Derdzi\'nski's Einstein classification
\cite[Proposition 5(i),(iv)]{Derdzinski1983} would make it nowhere
zero and would make $|W^+|^{2/3}g$ locally K\"ahler.  On the
orientation cover of its simple eigenline, of degree at most two,
the local K\"ahler forms glue: their only ambiguity is sign.
This gives a global K\"ahler metric on a finite cover, a
contradiction.  The same argument applies to $W^-$.

For a trace-free self-adjoint endomorphism $A$ of a three-dimensional
Euclidean space, its discriminant is
\[
 \mathfrak D(A)=\tfrac12\bigl(\operatorname{tr}A^2\bigr)^3
                  -3\bigl(\operatorname{tr}A^3\bigr)^2
              =\prod_{i<j}(\lambda_i-\lambda_j)^2.
\]
Einstein metrics are real analytic in harmonic coordinates
\cite[Theorem 5.2]{DeTurckKazdan}.  Thus
$\mathfrak D(W^\pm)$ are real-analytic functions, neither identically
zero.  Each zero set has empty interior, and the intersection of
their open dense complements proves the last assertion.
\end{proof}

These restrictions do not turn nonattainment of one curvature
functional into nonexistence of Einstein metrics.  Indeed, with
the full covariant tensor norm on $\Rm$,
\begin{equation}\label{eq:fixed-bundle-full-energy}
 \int_{\widehat E}|\Rm_g|^2\,d\mu
       =32\pi^2\widehat c+4\int_{\widehat E}|\Ric_g^\circ|^2\,d\mu
       =8V_*+4P(g).
\end{equation}
The explicit family makes $P\to0$, so the full-curvature infimum
is $8V_*$.  Its attainment is equivalent to existence of an Einstein
metric on $\widehat E$, because $\Ric^\circ=0$ and the contracted
Bianchi identity imply constant scalar curvature.  Such a metric
would have volume strictly greater than $V_*$; it need not attain
the scalar, Ricci, or constrained-volume infima.  Their values and
their different equality conditions are treated in
Section~\ref{sec:shear-variational}.

\subsection{Two precise restrictions on correction}

\begin{proposition}\label{prop:fixed-bundle-conformal-exclusion}
For every $\ell$, the conformal class of $G_\ell$ contains no
Einstein metric of any sign.  The assertion holds for its pullback
to every connected finite cover.
\end{proposition}

\begin{proof}
There is a nonempty open product region on which $G_\ell$ is the
product of two curvature-$-1$ surfaces.  Suppose
$\widehat G=f^{-2}G_\ell$ were Einstein, with $f>0$.  On a small
product neighborhood the trace-free conformal Ricci formula gives
\[
 0=\Ric_{\widehat G}^{\circ}
   =2f^{-1}(\nabla^2f)^\circ,
 \qquad \nabla^2f=\psi G_\ell.
\]
The mixed components imply $f=f_1(x)+f_2(y)$.  The equation on the
first factor makes $\psi$ independent of $y$, and that on the
second makes it independent of $x$; hence $\psi$ is constant.
Taking divergence on either factor, whose Ricci tensor is the
negative metric, gives $df_1=df_2=0$.  Thus $f$ is constant on
this product neighborhood.  The Einstein metric $\widehat G$ is
there locally symmetric, with nonpositive sectional curvature.
Its analyticity \cite{DeTurckKazdan} extends
$\widehat\nabla\Rm_{\widehat G}=0$ over the connected manifold.
Parallel transport then preserves the curvature tensor and
propagates nonpositive sectional curvature everywhere, contradicting
Proposition~\ref{prop:fixed-bundle-cover-obstructions}.  The
pulled-back family retains the same product neighborhoods, so the
argument is unchanged on a cover.
\end{proof}

\begin{proposition}\label{prop:fixed-bundle-path-exclusion}
The volume is constant along every $C^1$ path of smooth metrics
with $\Ric_g=-g$ on a fixed closed four-manifold.  Consequently
there is no $C^1$ family $\widehat G_\ell$, defined for all
$0<\ell<\ell_0$ on $E_D$, such that
\[
 \Ric_{\widehat G_\ell}=-\widehat G_\ell,\qquad
 \bigl\|G_\ell^{-1}(\widehat G_\ell-G_\ell)\bigr\|_\infty
       \longrightarrow0.
\]
The same statement holds on every fixed connected finite cover.
\end{proposition}

\begin{proof}
For a path $g(a)$ put $h=\partial_ag$ and $V(a)=\Vol(g(a))$.
The first variations of volume and total scalar curvature give
\[
 V'=\frac12\int\operatorname{tr}_gh\,d\mu,\qquad
 -4V'=\frac d{da}\int s_g\,d\mu
 =\int\left(-\langle\Ric_g,h\rangle
           +\frac{s_g}{2}\operatorname{tr}_gh\right)d\mu
 =-2V'.
\]
Hence $V'=0$.  Relative $C^0$ closeness to the explicit family
would make the constant volume equal to its limit $V_0$, contrary
to \eqref{eq:fixed-bundle-einstein-window}; use $V_*$ on a cover.
\end{proof}

The last two propositions rule out a conformal correction and a
nearby connected normalized Einstein branch, respectively.  They
do not rule out remote Einstein metrics or a discrete sequence of
exact metrics lying on different branches.  In particular, the
existence of any negative Einstein metric on these fixed smooth
manifolds remains a separate question.
